% Options for packages loaded elsewhere
\PassOptionsToPackage{unicode}{hyperref}
\PassOptionsToPackage{hyphens}{url}
\documentclass[
]{article}
\usepackage{xcolor}
\usepackage[margin=1in]{geometry}
\usepackage{amsmath,amssymb}
\setcounter{secnumdepth}{-\maxdimen} % remove section numbering
\usepackage{iftex}
\ifPDFTeX
  \usepackage[T1]{fontenc}
  \usepackage[utf8]{inputenc}
  \usepackage{textcomp} % provide euro and other symbols
\else % if luatex or xetex
  \usepackage{unicode-math} % this also loads fontspec
  \defaultfontfeatures{Scale=MatchLowercase}
  \defaultfontfeatures[\rmfamily]{Ligatures=TeX,Scale=1}
\fi
\usepackage{lmodern}
\ifPDFTeX\else
  % xetex/luatex font selection
\fi
% Use upquote if available, for straight quotes in verbatim environments
\IfFileExists{upquote.sty}{\usepackage{upquote}}{}
\IfFileExists{microtype.sty}{% use microtype if available
  \usepackage[]{microtype}
  \UseMicrotypeSet[protrusion]{basicmath} % disable protrusion for tt fonts
}{}
\makeatletter
\@ifundefined{KOMAClassName}{% if non-KOMA class
  \IfFileExists{parskip.sty}{%
    \usepackage{parskip}
  }{% else
    \setlength{\parindent}{0pt}
    \setlength{\parskip}{6pt plus 2pt minus 1pt}}
}{% if KOMA class
  \KOMAoptions{parskip=half}}
\makeatother
\usepackage{color}
\usepackage{fancyvrb}
\newcommand{\VerbBar}{|}
\newcommand{\VERB}{\Verb[commandchars=\\\{\}]}
\DefineVerbatimEnvironment{Highlighting}{Verbatim}{commandchars=\\\{\}}
% Add ',fontsize=\small' for more characters per line
\usepackage{framed}
\definecolor{shadecolor}{RGB}{248,248,248}
\newenvironment{Shaded}{\begin{snugshade}}{\end{snugshade}}
\newcommand{\AlertTok}[1]{\textcolor[rgb]{0.94,0.16,0.16}{#1}}
\newcommand{\AnnotationTok}[1]{\textcolor[rgb]{0.56,0.35,0.01}{\textbf{\textit{#1}}}}
\newcommand{\AttributeTok}[1]{\textcolor[rgb]{0.13,0.29,0.53}{#1}}
\newcommand{\BaseNTok}[1]{\textcolor[rgb]{0.00,0.00,0.81}{#1}}
\newcommand{\BuiltInTok}[1]{#1}
\newcommand{\CharTok}[1]{\textcolor[rgb]{0.31,0.60,0.02}{#1}}
\newcommand{\CommentTok}[1]{\textcolor[rgb]{0.56,0.35,0.01}{\textit{#1}}}
\newcommand{\CommentVarTok}[1]{\textcolor[rgb]{0.56,0.35,0.01}{\textbf{\textit{#1}}}}
\newcommand{\ConstantTok}[1]{\textcolor[rgb]{0.56,0.35,0.01}{#1}}
\newcommand{\ControlFlowTok}[1]{\textcolor[rgb]{0.13,0.29,0.53}{\textbf{#1}}}
\newcommand{\DataTypeTok}[1]{\textcolor[rgb]{0.13,0.29,0.53}{#1}}
\newcommand{\DecValTok}[1]{\textcolor[rgb]{0.00,0.00,0.81}{#1}}
\newcommand{\DocumentationTok}[1]{\textcolor[rgb]{0.56,0.35,0.01}{\textbf{\textit{#1}}}}
\newcommand{\ErrorTok}[1]{\textcolor[rgb]{0.64,0.00,0.00}{\textbf{#1}}}
\newcommand{\ExtensionTok}[1]{#1}
\newcommand{\FloatTok}[1]{\textcolor[rgb]{0.00,0.00,0.81}{#1}}
\newcommand{\FunctionTok}[1]{\textcolor[rgb]{0.13,0.29,0.53}{\textbf{#1}}}
\newcommand{\ImportTok}[1]{#1}
\newcommand{\InformationTok}[1]{\textcolor[rgb]{0.56,0.35,0.01}{\textbf{\textit{#1}}}}
\newcommand{\KeywordTok}[1]{\textcolor[rgb]{0.13,0.29,0.53}{\textbf{#1}}}
\newcommand{\NormalTok}[1]{#1}
\newcommand{\OperatorTok}[1]{\textcolor[rgb]{0.81,0.36,0.00}{\textbf{#1}}}
\newcommand{\OtherTok}[1]{\textcolor[rgb]{0.56,0.35,0.01}{#1}}
\newcommand{\PreprocessorTok}[1]{\textcolor[rgb]{0.56,0.35,0.01}{\textit{#1}}}
\newcommand{\RegionMarkerTok}[1]{#1}
\newcommand{\SpecialCharTok}[1]{\textcolor[rgb]{0.81,0.36,0.00}{\textbf{#1}}}
\newcommand{\SpecialStringTok}[1]{\textcolor[rgb]{0.31,0.60,0.02}{#1}}
\newcommand{\StringTok}[1]{\textcolor[rgb]{0.31,0.60,0.02}{#1}}
\newcommand{\VariableTok}[1]{\textcolor[rgb]{0.00,0.00,0.00}{#1}}
\newcommand{\VerbatimStringTok}[1]{\textcolor[rgb]{0.31,0.60,0.02}{#1}}
\newcommand{\WarningTok}[1]{\textcolor[rgb]{0.56,0.35,0.01}{\textbf{\textit{#1}}}}
\usepackage{graphicx}
\makeatletter
\newsavebox\pandoc@box
\newcommand*\pandocbounded[1]{% scales image to fit in text height/width
  \sbox\pandoc@box{#1}%
  \Gscale@div\@tempa{\textheight}{\dimexpr\ht\pandoc@box+\dp\pandoc@box\relax}%
  \Gscale@div\@tempb{\linewidth}{\wd\pandoc@box}%
  \ifdim\@tempb\p@<\@tempa\p@\let\@tempa\@tempb\fi% select the smaller of both
  \ifdim\@tempa\p@<\p@\scalebox{\@tempa}{\usebox\pandoc@box}%
  \else\usebox{\pandoc@box}%
  \fi%
}
% Set default figure placement to htbp
\def\fps@figure{htbp}
\makeatother
\setlength{\emergencystretch}{3em} % prevent overfull lines
\providecommand{\tightlist}{%
  \setlength{\itemsep}{0pt}\setlength{\parskip}{0pt}}
\usepackage{bookmark}
\IfFileExists{xurl.sty}{\usepackage{xurl}}{} % add URL line breaks if available
\urlstyle{same}
\hypersetup{
  pdftitle={Introduction to NonlinearBSS},
  hidelinks,
  pdfcreator={LaTeX via pandoc}}

\title{Introduction to NonlinearBSS}
\author{}
\date{\vspace{-2.5em}}

\begin{document}
\maketitle

In nonlinear blind source separation the observed \(p\)-variate data
\(x\) are assumed to be generated as \(x = f(z)\), where \(f\) is an
unknown injective function and \(z\) are independent latent components.
The aim is to recover \(z\) from \(x\). Without further assumptions this
is impossible, but nonstationarity of the components over time or space,
or temporal dependence, makes them identifiable.

This vignette illustrates three methods:

\begin{itemize}
\tightlist
\item
  \texttt{VAE()}: a variational autoencoder for i.i.d. data (a baseline;
  not identifiable),
\item
  \texttt{iVAE()}: the identifiable variational autoencoder for any data
  with auxiliary variables,
\item
  \texttt{iVAEar()}: iVAE with nonstationary autoregressive latent
  components.
\end{itemize}

The graph-based \texttt{GiVAE()} for areal data and networks is
described in its own vignettes. All plots of the package are made with
\textbf{ggplot2} and return \texttt{ggplot} objects that can be modified
further.

The models are fitted with TensorFlow through \textbf{keras3}. Install
it once with \texttt{keras3::install\_keras()} and check the setup with
\texttt{tf\_available()}.

\subsection{Simulated data}\label{simulated-data}

\texttt{simulate\_sources()} creates components whose means and
variances change over segments and \texttt{mix\_data()} mixes them with
a random neural network.

\begin{Shaded}
\begin{Highlighting}[]
\FunctionTok{set.seed}\NormalTok{(}\DecValTok{1}\NormalTok{)}
\NormalTok{src }\OtherTok{\textless{}{-}} \FunctionTok{simulate\_sources}\NormalTok{(}\DecValTok{3}\NormalTok{, }\AttributeTok{time =} \DecValTok{1}\SpecialCharTok{:}\DecValTok{3000}\NormalTok{, }\AttributeTok{n\_segments\_time =} \DecValTok{30}\NormalTok{)}
\NormalTok{x }\OtherTok{\textless{}{-}} \FunctionTok{mix\_data}\NormalTok{(src}\SpecialCharTok{$}\NormalTok{sources, }\AttributeTok{n\_layers =} \DecValTok{3}\NormalTok{)}
\end{Highlighting}
\end{Shaded}

\subsection{Auxiliary data}\label{auxiliary-data}

For iVAE the auxiliary variables \(u\) should explain how the
distribution of the components changes. For a time series, indicators of
time segments are a simple choice:

\begin{Shaded}
\begin{Highlighting}[]
\NormalTok{A }\OtherTok{\textless{}{-}} \FunctionTok{aux\_matrix}\NormalTok{(}\FunctionTok{aux\_segments}\NormalTok{(}\DecValTok{100}\NormalTok{), x, }\AttributeTok{time =} \DecValTok{1}\SpecialCharTok{:}\DecValTok{3000}\NormalTok{)}
\FunctionTok{dim}\NormalTok{(A)}
\CommentTok{\#\textgreater{} [1] 3000   30}
\end{Highlighting}
\end{Shaded}

Radial basis functions give a smooth and more flexible alternative:

\begin{Shaded}
\begin{Highlighting}[]
\NormalTok{A }\OtherTok{\textless{}{-}} \FunctionTok{aux\_matrix}\NormalTok{(}\FunctionTok{aux\_rbf}\NormalTok{(}\AttributeTok{temporal\_basis =} \FunctionTok{c}\NormalTok{(}\DecValTok{9}\NormalTok{, }\DecValTok{17}\NormalTok{, }\DecValTok{37}\NormalTok{)), x, }\AttributeTok{time =} \DecValTok{1}\SpecialCharTok{:}\DecValTok{3000}\NormalTok{)}
\FunctionTok{dim}\NormalTok{(A)}
\CommentTok{\#\textgreater{} [1] 3000   63}
\end{Highlighting}
\end{Shaded}

\subsection{Fitting iVAE}\label{fitting-ivae}

\begin{Shaded}
\begin{Highlighting}[]
\NormalTok{fit }\OtherTok{\textless{}{-}} \FunctionTok{iVAE}\NormalTok{(x, }\FunctionTok{aux\_rbf}\NormalTok{(}\AttributeTok{temporal\_basis =} \FunctionTok{c}\NormalTok{(}\DecValTok{9}\NormalTok{, }\DecValTok{17}\NormalTok{, }\DecValTok{37}\NormalTok{)), }\AttributeTok{time =} \DecValTok{1}\SpecialCharTok{:}\DecValTok{3000}\NormalTok{, }\AttributeTok{epochs =} \DecValTok{50}\NormalTok{,}
            \AttributeTok{verbose =} \DecValTok{0}\NormalTok{, }\AttributeTok{seed =} \DecValTok{1}\NormalTok{)}
\NormalTok{fit}
\CommentTok{\#\textgreater{} Identifiable variational autoencoder (iVAE)}
\CommentTok{\#\textgreater{} }
\CommentTok{\#\textgreater{} Call: iVAE(x = x, aux = aux\_rbf(temporal\_basis = c(9, 17, 37)), time = 1:3000, }
\CommentTok{\#\textgreater{}     epochs = 50, seed = 1, verbose = 0)}
\CommentTok{\#\textgreater{} Data: x (3000 observations, 3 variables, format matrix)}
\CommentTok{\#\textgreater{} Latent components: 3}
\CommentTok{\#\textgreater{} Auxiliary data: radial basis functions, spatial = 2, 9, temporal = 9, 17, 37 (gaussian) [63 variables]}
\CommentTok{\#\textgreater{} Distributions: source gaussian, error gaussian (sigma 0.02)}
\CommentTok{\#\textgreater{} ELBO: {-}3.1172}
\FunctionTok{absolute\_mean\_correlation}\NormalTok{(fit}\SpecialCharTok{$}\NormalTok{S, src}\SpecialCharTok{$}\NormalTok{sources)}
\CommentTok{\#\textgreater{} [1] 0.9565594}
\end{Highlighting}
\end{Shaded}

For comparison, the plain \texttt{VAE()} has no auxiliary variables. It
is a baseline whose components are not identifiable in general, so its
recovery is not guaranteed:

\begin{Shaded}
\begin{Highlighting}[]
\NormalTok{fit\_vae }\OtherTok{\textless{}{-}} \FunctionTok{VAE}\NormalTok{(x, }\AttributeTok{latent\_dim =} \DecValTok{3}\NormalTok{, }\AttributeTok{epochs =} \DecValTok{50}\NormalTok{, }\AttributeTok{verbose =} \DecValTok{0}\NormalTok{, }\AttributeTok{seed =} \DecValTok{1}\NormalTok{)}
\FunctionTok{absolute\_mean\_correlation}\NormalTok{(fit\_vae}\SpecialCharTok{$}\NormalTok{S, src}\SpecialCharTok{$}\NormalTok{sources)}
\CommentTok{\#\textgreater{} [1] 0.7132899}
\end{Highlighting}
\end{Shaded}

The mean correlation coefficient compares the estimates with the true
components up to permutation and sign; 1 means perfect recovery. The
components are stored as \texttt{fit\$S} and are returned in the format
of the data by \texttt{components()}.

\begin{Shaded}
\begin{Highlighting}[]
\FunctionTok{head}\NormalTok{(}\FunctionTok{components}\NormalTok{(fit))}
\CommentTok{\#\textgreater{}             IC1        IC2        IC3}
\CommentTok{\#\textgreater{} [1,]  2.0283578  0.1393474 {-}0.4877269}
\CommentTok{\#\textgreater{} [2,] {-}1.6412711 {-}0.6063424 {-}1.6819368}
\CommentTok{\#\textgreater{} [3,]  1.0652925 {-}0.6113633 {-}1.2268908}
\CommentTok{\#\textgreater{} [4,] {-}0.7140939  0.9349620 {-}0.6195569}
\CommentTok{\#\textgreater{} [5,]  0.4349818 {-}1.6999949 {-}1.6443090}
\CommentTok{\#\textgreater{} [6,] {-}0.8067482  0.4258073 {-}0.2712600}
\end{Highlighting}
\end{Shaded}

\subsection{Choosing the number of latent
components}\label{choosing-the-number-of-latent-components}

In practice the number of latent components is unknown. Following Sipilä
et al.~(2025, \emph{Neural Networks}; 2026, \emph{Spatial Statistics}),
it can be selected by fitting models with different numbers of
components \(L\) and maximising the profile AIC
\[\mathrm{pAIC} = \mathrm{ELBO} - L,\] where the ELBO is the average
evidence lower bound per observation. Here the three sources are mixed
into six observed variables, and models with one to six components are
compared:

\begin{Shaded}
\begin{Highlighting}[]
\NormalTok{x6 }\OtherTok{\textless{}{-}} \FunctionTok{mix\_data}\NormalTok{(src}\SpecialCharTok{$}\NormalTok{sources, }\AttributeTok{n\_layers =} \DecValTok{3}\NormalTok{, }\AttributeTok{output\_dim =} \DecValTok{6}\NormalTok{)}
\NormalTok{sel }\OtherTok{\textless{}{-}} \FunctionTok{select\_latent\_dim}\NormalTok{(x6, }\AttributeTok{latent\_dims =} \DecValTok{1}\SpecialCharTok{:}\DecValTok{6}\NormalTok{, }\AttributeTok{aux =} \FunctionTok{aux\_segments}\NormalTok{(}\DecValTok{100}\NormalTok{),}
                         \AttributeTok{time =} \DecValTok{1}\SpecialCharTok{:}\DecValTok{3000}\NormalTok{, }\AttributeTok{epochs =} \DecValTok{50}\NormalTok{, }\AttributeTok{verbose =} \DecValTok{0}\NormalTok{, }\AttributeTok{seed =} \DecValTok{1}\NormalTok{)}
\NormalTok{sel}
\CommentTok{\#\textgreater{} Selection of the number of latent components}
\CommentTok{\#\textgreater{} }
\CommentTok{\#\textgreater{}  latent\_dim    elbo    pAIC selected}
\CommentTok{\#\textgreater{}           1 {-}31.703 {-}32.703         }
\CommentTok{\#\textgreater{}           2  {-}8.562 {-}10.562         }
\CommentTok{\#\textgreater{}           3  {-}1.491  {-}4.491        *}
\CommentTok{\#\textgreater{}           4  {-}1.696  {-}5.696         }
\CommentTok{\#\textgreater{}           5  {-}1.585  {-}6.585         }
\CommentTok{\#\textgreater{}           6  {-}1.631  {-}7.631         }
\CommentTok{\#\textgreater{} }
\CommentTok{\#\textgreater{} Selected by the largest pAIC = ELBO {-} number of components: 3}
\end{Highlighting}
\end{Shaded}

The ELBO increases steeply until the true number of components and stays
nearly constant after that:

\begin{Shaded}
\begin{Highlighting}[]
\FunctionTok{plot}\NormalTok{(sel)}
\end{Highlighting}
\end{Shaded}

\includegraphics[width=1\linewidth]{figures/dimsel-1}

The selected model is kept in the result. Models that were fitted
separately, for example in different sessions, can be compared with
\texttt{latent\_dim\_selection(list(fit1,\ fit2,\ ...))}. The models
should differ only in the number of components, since the ELBO also
depends on the data, the auxiliary variables and
\texttt{error\_dist\_sigma}.

\subsection{Time series formats}\label{time-series-formats}

Time series objects (\texttt{ts}, \texttt{xts}, \texttt{zoo},
\texttt{tsibble}) can be used directly, and the results are returned in
the same format.

The time points of these objects are measured in observation steps (here
months), as for matrices, so the auxiliary variables are specified in
the same way for all data formats:

\begin{Shaded}
\begin{Highlighting}[]
\NormalTok{xt }\OtherTok{\textless{}{-}} \FunctionTok{ts}\NormalTok{(x, }\AttributeTok{start =} \FunctionTok{c}\NormalTok{(}\DecValTok{1800}\NormalTok{, }\DecValTok{1}\NormalTok{), }\AttributeTok{frequency =} \DecValTok{12}\NormalTok{)}
\NormalTok{fit\_ts }\OtherTok{\textless{}{-}} \FunctionTok{iVAE}\NormalTok{(xt, }\FunctionTok{aux\_segments}\NormalTok{(}\DecValTok{100}\NormalTok{), }\AttributeTok{epochs =} \DecValTok{50}\NormalTok{, }\AttributeTok{verbose =} \DecValTok{0}\NormalTok{, }\AttributeTok{seed =} \DecValTok{1}\NormalTok{)}
\FunctionTok{class}\NormalTok{(}\FunctionTok{components}\NormalTok{(fit\_ts))}
\CommentTok{\#\textgreater{} [1] "mts"    "ts"     "matrix" "array"}
\FunctionTok{absolute\_mean\_correlation}\NormalTok{(fit\_ts}\SpecialCharTok{$}\NormalTok{S, src}\SpecialCharTok{$}\NormalTok{sources)}
\CommentTok{\#\textgreater{} [1] 0.7475805}
\end{Highlighting}
\end{Shaded}

\begin{Shaded}
\begin{Highlighting}[]
\FunctionTok{plot}\NormalTok{(fit\_ts)}
\end{Highlighting}
\end{Shaded}

\includegraphics[width=1\linewidth]{figures/ts-components-1}

\subsection{A second simulated example: spatio-temporal
data}\label{a-second-simulated-example-spatio-temporal-data}

\texttt{simulate\_sources()} also creates spatio-temporal components,
with nonstationarity over space and time and, optionally, autoregressive
dependence. Here three latent components are observed at 60 random
locations and at 40 time points, and the three components are mixed into
five observed variables with a two-layer random mixing function:

\begin{Shaded}
\begin{Highlighting}[]
\FunctionTok{set.seed}\NormalTok{(}\DecValTok{1}\NormalTok{)}
\NormalTok{n\_s }\OtherTok{\textless{}{-}} \DecValTok{60}
\NormalTok{n\_t }\OtherTok{\textless{}{-}} \DecValTok{100}
\NormalTok{xy }\OtherTok{\textless{}{-}} \FunctionTok{cbind}\NormalTok{(}\AttributeTok{lon =} \FunctionTok{runif}\NormalTok{(n\_s), }\AttributeTok{lat =} \FunctionTok{runif}\NormalTok{(n\_s))}
\NormalTok{coords }\OtherTok{\textless{}{-}}\NormalTok{ xy[}\FunctionTok{rep}\NormalTok{(}\FunctionTok{seq\_len}\NormalTok{(n\_s), n\_t), ]}
\NormalTok{time }\OtherTok{\textless{}{-}} \FunctionTok{rep}\NormalTok{(}\FunctionTok{seq\_len}\NormalTok{(n\_t), }\AttributeTok{each =}\NormalTok{ n\_s)}
\NormalTok{src\_st }\OtherTok{\textless{}{-}} \FunctionTok{simulate\_sources}\NormalTok{(}\DecValTok{3}\NormalTok{, }\AttributeTok{coords =}\NormalTok{ coords, }\AttributeTok{time =}\NormalTok{ time,}
                           \AttributeTok{n\_segments\_space =} \DecValTok{6}\NormalTok{, }\AttributeTok{n\_segments\_time =} \DecValTok{5}\NormalTok{,}
                           \AttributeTok{ar\_range =} \FunctionTok{c}\NormalTok{(}\FloatTok{0.3}\NormalTok{, }\FloatTok{0.9}\NormalTok{))}
\NormalTok{x\_st }\OtherTok{\textless{}{-}} \FunctionTok{mix\_data}\NormalTok{(src\_st}\SpecialCharTok{$}\NormalTok{sources, }\AttributeTok{n\_layers =} \DecValTok{2}\NormalTok{, }\AttributeTok{output\_dim =} \DecValTok{5}\NormalTok{)}
\NormalTok{dat }\OtherTok{\textless{}{-}} \FunctionTok{data.frame}\NormalTok{(x\_st, coords, }\AttributeTok{time =}\NormalTok{ time)}
\NormalTok{vars\_st }\OtherTok{\textless{}{-}} \FunctionTok{colnames}\NormalTok{(x\_st)}
\end{Highlighting}
\end{Shaded}

The data frame is used with the names of the coordinate, time and
variable columns. The radial basis functions of \texttt{aux\_rbf()} are
the auxiliary variables. \texttt{VAE()} ignores the coordinates and
times, \texttt{iVAE()} treats the observations as independent given the
auxiliary variables, and \texttt{iVAEar()} also models the
autoregressive dependence:

\begin{Shaded}
\begin{Highlighting}[]
\NormalTok{aux }\OtherTok{\textless{}{-}} \FunctionTok{aux\_rbf}\NormalTok{(}\AttributeTok{spatial\_basis =} \FunctionTok{c}\NormalTok{(}\DecValTok{2}\NormalTok{, }\DecValTok{9}\NormalTok{), }\AttributeTok{temporal\_basis =} \FunctionTok{c}\NormalTok{(}\DecValTok{9}\NormalTok{, }\DecValTok{17}\NormalTok{))}
\NormalTok{fit\_st\_vae }\OtherTok{\textless{}{-}} \FunctionTok{VAE}\NormalTok{(dat[vars\_st], }\AttributeTok{latent\_dim =} \DecValTok{3}\NormalTok{, }\AttributeTok{epochs =} \DecValTok{60}\NormalTok{, }\AttributeTok{verbose =} \DecValTok{0}\NormalTok{, }\AttributeTok{seed =} \DecValTok{1}\NormalTok{)}
\NormalTok{fit\_st\_ivae }\OtherTok{\textless{}{-}} \FunctionTok{iVAE}\NormalTok{(dat, aux, }\AttributeTok{coords =} \FunctionTok{c}\NormalTok{(}\StringTok{"lon"}\NormalTok{, }\StringTok{"lat"}\NormalTok{), }\AttributeTok{time =} \StringTok{"time"}\NormalTok{,}
                    \AttributeTok{vars =}\NormalTok{ vars\_st, }\AttributeTok{latent\_dim =} \DecValTok{3}\NormalTok{, }\AttributeTok{epochs =} \DecValTok{60}\NormalTok{, }\AttributeTok{verbose =} \DecValTok{0}\NormalTok{,}
                    \AttributeTok{seed =} \DecValTok{1}\NormalTok{)}
\NormalTok{fit\_st\_ar }\OtherTok{\textless{}{-}} \FunctionTok{iVAEar}\NormalTok{(dat, aux, }\AttributeTok{coords =} \FunctionTok{c}\NormalTok{(}\StringTok{"lon"}\NormalTok{, }\StringTok{"lat"}\NormalTok{), }\AttributeTok{time =} \StringTok{"time"}\NormalTok{,}
                    \AttributeTok{vars =}\NormalTok{ vars\_st, }\AttributeTok{latent\_dim =} \DecValTok{3}\NormalTok{, }\AttributeTok{epochs =} \DecValTok{60}\NormalTok{, }\AttributeTok{verbose =} \DecValTok{0}\NormalTok{,}
                    \AttributeTok{seed =} \DecValTok{1}\NormalTok{)}
\FunctionTok{sapply}\NormalTok{(}\FunctionTok{list}\NormalTok{(}\AttributeTok{VAE =}\NormalTok{ fit\_st\_vae, }\AttributeTok{iVAE =}\NormalTok{ fit\_st\_ivae, }\AttributeTok{iVAEar =}\NormalTok{ fit\_st\_ar),}
       \ControlFlowTok{function}\NormalTok{(f) }\FunctionTok{absolute\_mean\_correlation}\NormalTok{(f}\SpecialCharTok{$}\NormalTok{S, src\_st}\SpecialCharTok{$}\NormalTok{sources))}
\CommentTok{\#\textgreater{}       VAE      iVAE    iVAEar }
\CommentTok{\#\textgreater{} 0.8760115 0.9822837 0.9917308}
\end{Highlighting}
\end{Shaded}

The estimated components as a map at the last time point, and as a time
series at the first location:

\begin{Shaded}
\begin{Highlighting}[]
\FunctionTok{plot}\NormalTok{(fit\_st\_ar)}
\end{Highlighting}
\end{Shaded}

\includegraphics[width=1\linewidth]{figures/st-map-1}

\begin{Shaded}
\begin{Highlighting}[]
\FunctionTok{plot}\NormalTok{(fit\_st\_ar, }\AttributeTok{type =} \StringTok{"series"}\NormalTok{, }\AttributeTok{series =} \DecValTok{1}\NormalTok{)}
\end{Highlighting}
\end{Shaded}

\includegraphics[width=1\linewidth]{figures/st-series-1}

Since the result is a \texttt{ggplot} object, it can be customised with
the usual ggplot2 functions, for example
\texttt{plot(fit\_st\_ar)\ +\ ggplot2::ggtitle("iVAEar")}.

\subsection{Spatio-temporal data}\label{spatio-temporal-data}

The data set \texttt{pollution\_northern\_italy} contains daily air
quality measurements at 50 stations. Here, the observations from January
2020 to April 2022 are used for training, May 2022 as forecasting test
data. Five stations are set aside to test spatial interpolation. The
data frame is used with the column names of the coordinates, time points
and variables.

\begin{Shaded}
\begin{Highlighting}[]
\NormalTok{train\_subset\_time }\OtherTok{\textless{}{-}} \FunctionTok{subset}\NormalTok{(pollution\_northern\_italy, time }\SpecialCharTok{\textless{}} \FunctionTok{as.Date}\NormalTok{(}\StringTok{"2022{-}05{-}01"}\NormalTok{))}
\NormalTok{stations }\OtherTok{\textless{}{-}} \FunctionTok{unique}\NormalTok{(train\_subset\_time}\SpecialCharTok{$}\NormalTok{AirQualityStation)}
\FunctionTok{set.seed}\NormalTok{(}\DecValTok{2}\NormalTok{)}
\NormalTok{test\_stations }\OtherTok{\textless{}{-}} \FunctionTok{sample}\NormalTok{(stations, }\DecValTok{5}\NormalTok{)}
\NormalTok{train\_subset }\OtherTok{\textless{}{-}} \FunctionTok{subset}\NormalTok{(train\_subset\_time, }\SpecialCharTok{!}\NormalTok{AirQualityStation }\SpecialCharTok{\%in\%}\NormalTok{ test\_stations)}
\NormalTok{vars }\OtherTok{\textless{}{-}} \FunctionTok{c}\NormalTok{(}\StringTok{"mean\_O3"}\NormalTok{, }\StringTok{"mean\_NO2"}\NormalTok{, }\StringTok{"mean\_PM10"}\NormalTok{)}
\end{Highlighting}
\end{Shaded}

\texttt{iVAEar()} models the components as nonstationary autoregressive
processes. The model is trained for 30 epochs with RBF auxiliary
variables. The temporal RBFs are formed based on the seasonal time point
(1 to 365) instead of the absolute time point in order to allow the
model to generalize better to future.

\begin{Shaded}
\begin{Highlighting}[]
\NormalTok{fit }\OtherTok{\textless{}{-}} \FunctionTok{iVAEar}\NormalTok{(train\_subset, }\AttributeTok{aux =} \FunctionTok{aux\_rbf}\NormalTok{(}\AttributeTok{spatial\_basis =} \FunctionTok{c}\NormalTok{(}\DecValTok{2}\NormalTok{, }\DecValTok{9}\NormalTok{),}
                                   \AttributeTok{temporal\_basis =} \FunctionTok{c}\NormalTok{(}\DecValTok{9}\NormalTok{, }\DecValTok{17}\NormalTok{), }\AttributeTok{seasonal\_period =} \DecValTok{365}\NormalTok{),}
              \AttributeTok{coords =} \FunctionTok{c}\NormalTok{(}\StringTok{"Longitude"}\NormalTok{, }\StringTok{"Latitude"}\NormalTok{), }\AttributeTok{time =} \StringTok{"time"}\NormalTok{, }\AttributeTok{vars =}\NormalTok{ vars,}
              \AttributeTok{epochs =} \DecValTok{30}\NormalTok{, }\AttributeTok{verbose =} \DecValTok{0}\NormalTok{, }\AttributeTok{seed =} \DecValTok{1}\NormalTok{)}
\FunctionTok{summary}\NormalTok{(fit)}
\CommentTok{\#\textgreater{} Identifiable variational autoencoder with autoregressive components (iVAEar)}
\CommentTok{\#\textgreater{} }
\CommentTok{\#\textgreater{} Call: iVAEar(x = train\_subset, aux = aux\_rbf(spatial\_basis = c(2, 9), }
\CommentTok{\#\textgreater{}     temporal\_basis = c(9, 17), seasonal\_period = 365), coords = c("Longitude", }
\CommentTok{\#\textgreater{}     "Latitude"), time = "time", vars = vars, epochs = 30, seed = 1, }
\CommentTok{\#\textgreater{}     verbose = 0)}
\CommentTok{\#\textgreater{} }
\CommentTok{\#\textgreater{} Data: train\_subset }
\CommentTok{\#\textgreater{}   observations: 38295 }
\CommentTok{\#\textgreater{}   variables: 3 }
\CommentTok{\#\textgreater{}   format: data.frame }
\CommentTok{\#\textgreater{} Latent components: 3 }
\CommentTok{\#\textgreater{} Autoregressive order: 1 }
\CommentTok{\#\textgreater{} Auxiliary data: radial basis functions, spatial = 2, 9, temporal = 9, 17 (gaussian) + seasonal period 365 [114 variables]}
\CommentTok{\#\textgreater{} }
\CommentTok{\#\textgreater{} Network: hidden units 128{-}128{-}128, auxiliary function 128{-}128{-}128, activation leaky\_relu}
\CommentTok{\#\textgreater{} Distributions: source gaussian, error gaussian}
\CommentTok{\#\textgreater{} }
\CommentTok{\#\textgreater{} Training: 30 epochs, final loss 0.975}
\CommentTok{\#\textgreater{} ELBO: {-}0.95145 }
\CommentTok{\#\textgreater{} Largest absolute correlation between components: 0.463 }
\CommentTok{\#\textgreater{} }
\CommentTok{\#\textgreater{} Prior variance of the standardised components over the data:}
\CommentTok{\#\textgreater{}         0\%    50\%  100\%}
\CommentTok{\#\textgreater{} IC1 0.0118 0.0947 0.897}
\CommentTok{\#\textgreater{} IC2 0.0339 0.1442 0.690}
\CommentTok{\#\textgreater{} IC3 0.1039 0.3463 1.741}
\CommentTok{\#\textgreater{} }
\CommentTok{\#\textgreater{} Autoregressive coefficients over the data:}
\CommentTok{\#\textgreater{}  component lag     min  mean   max}
\CommentTok{\#\textgreater{}        IC1   1 {-}0.1478 0.566 0.980}
\CommentTok{\#\textgreater{}        IC2   1  0.1481 0.656 0.960}
\CommentTok{\#\textgreater{}        IC3   1  0.0871 0.579 0.855}
\end{Highlighting}
\end{Shaded}

\begin{Shaded}
\begin{Highlighting}[]
\FunctionTok{plot}\NormalTok{(fit, }\AttributeTok{type =} \StringTok{"history"}\NormalTok{)}
\end{Highlighting}
\end{Shaded}

\includegraphics[width=1\linewidth]{figures/history-1}

The components can be shown as maps at a given day, or as time series at
a station:

\begin{Shaded}
\begin{Highlighting}[]
\FunctionTok{plot}\NormalTok{(fit, }\AttributeTok{time =} \FunctionTok{as.Date}\NormalTok{(}\StringTok{"2022{-}04{-}01"}\NormalTok{))}
\end{Highlighting}
\end{Shaded}

\includegraphics[width=1\linewidth]{figures/maps-1}

\begin{Shaded}
\begin{Highlighting}[]
\FunctionTok{plot}\NormalTok{(fit, }\AttributeTok{type =} \StringTok{"series"}\NormalTok{, }\AttributeTok{series =} \DecValTok{1}\NormalTok{)}
\end{Highlighting}
\end{Shaded}

\includegraphics[width=1\linewidth]{figures/series-1}

\subsection{Forecasting and
interpolation}\label{forecasting-and-interpolation}

The trend, variance and autoregressive coefficient functions can be
evaluated anywhere, and the components can be forecast from their last
values by iterating the autoregressive prior. With
\texttt{space\ =\ "data"} the predicted components are decoded to the
observed variables. Here we forecast the first week of May 2022, one to
seven days after the end of the training data:

\begin{Shaded}
\begin{Highlighting}[]
\NormalTok{test\_time }\OtherTok{\textless{}{-}} \FunctionTok{subset}\NormalTok{(pollution\_northern\_italy, }
\NormalTok{                     time }\SpecialCharTok{\textless{}} \FunctionTok{as.Date}\NormalTok{(}\StringTok{"2022{-}05{-}08"}\NormalTok{) }\SpecialCharTok{\&}\NormalTok{ time }\SpecialCharTok{\textgreater{}} \FunctionTok{as.Date}\NormalTok{(}\StringTok{"2022{-}04{-}30"}\NormalTok{))}
\NormalTok{pred\_time }\OtherTok{\textless{}{-}} \FunctionTok{predict}\NormalTok{(fit, test\_time, }\AttributeTok{type =} \StringTok{"forecast"}\NormalTok{,}
              \AttributeTok{coords =} \FunctionTok{c}\NormalTok{(}\StringTok{"Longitude"}\NormalTok{, }\StringTok{"Latitude"}\NormalTok{), }\AttributeTok{time =} \StringTok{"time"}\NormalTok{,}
              \AttributeTok{space =} \StringTok{"data"}\NormalTok{)}
\FunctionTok{head}\NormalTok{(pred\_time)}
\CommentTok{\#\textgreater{}       Longitude Latitude       time   mean\_O3 mean\_NO2 mean\_PM10}
\CommentTok{\#\textgreater{} 63663  11.43317 46.58967 2022{-}05{-}01 104.59662 4.334665  12.58798}
\CommentTok{\#\textgreater{} 63664  11.43317 46.58967 2022{-}05{-}02  98.84446 4.166092  12.81774}
\CommentTok{\#\textgreater{} 63665  11.43317 46.58967 2022{-}05{-}03  94.61120 4.104034  12.39065}
\CommentTok{\#\textgreater{} 63666  11.43317 46.58967 2022{-}05{-}04  91.35274 4.034678  12.00895}
\CommentTok{\#\textgreater{} 63667  11.43317 46.58967 2022{-}05{-}05  88.59262 3.947743  11.64543}
\CommentTok{\#\textgreater{} 63668  11.43317 46.58967 2022{-}05{-}06  86.35072 3.904140  11.32179}
\end{Highlighting}
\end{Shaded}

The mean absolute error of the ozone forecasts by forecast day, compared
with the standard deviation of ozone in the test period:

\begin{Shaded}
\begin{Highlighting}[]
\FunctionTok{round}\NormalTok{(}\FunctionTok{c}\NormalTok{(}\AttributeTok{MAE =} \FunctionTok{mean}\NormalTok{(}\FunctionTok{abs}\NormalTok{(pred\_time}\SpecialCharTok{$}\NormalTok{mean\_O3 }\SpecialCharTok{{-}}\NormalTok{ test\_time}\SpecialCharTok{$}\NormalTok{mean\_O3)),}
        \AttributeTok{sd =} \FunctionTok{sd}\NormalTok{(test\_time}\SpecialCharTok{$}\NormalTok{mean\_O3)), }\DecValTok{1}\NormalTok{)}
\CommentTok{\#\textgreater{}  MAE   sd }
\CommentTok{\#\textgreater{} 13.5 18.8}
\end{Highlighting}
\end{Shaded}

Spatial interpolation uses the trend function at the locations of the
five stations not used in fitting:

\begin{Shaded}
\begin{Highlighting}[]
\NormalTok{test\_space }\OtherTok{\textless{}{-}} \FunctionTok{subset}\NormalTok{(train\_subset\_time, AirQualityStation }\SpecialCharTok{\%in\%}\NormalTok{ test\_stations)}
\NormalTok{pred\_space }\OtherTok{\textless{}{-}} \FunctionTok{predict}\NormalTok{(fit, test\_space, }\AttributeTok{type =} \StringTok{"trend"}\NormalTok{,}
                  \AttributeTok{coords =} \FunctionTok{c}\NormalTok{(}\StringTok{"Longitude"}\NormalTok{, }\StringTok{"Latitude"}\NormalTok{), }\AttributeTok{time =} \StringTok{"time"}\NormalTok{,}
                  \AttributeTok{space =} \StringTok{"data"}\NormalTok{)}
\FunctionTok{round}\NormalTok{(}\FunctionTok{c}\NormalTok{(}\AttributeTok{MAE =} \FunctionTok{mean}\NormalTok{(}\FunctionTok{abs}\NormalTok{(pred\_space}\SpecialCharTok{$}\NormalTok{mean\_O3 }\SpecialCharTok{{-}}\NormalTok{ test\_space}\SpecialCharTok{$}\NormalTok{mean\_O3)),}
        \AttributeTok{sd =} \FunctionTok{sd}\NormalTok{(test\_space}\SpecialCharTok{$}\NormalTok{mean\_O3)), }\DecValTok{1}\NormalTok{)}
\CommentTok{\#\textgreater{}  MAE   sd }
\CommentTok{\#\textgreater{} 14.9 33.2}
\end{Highlighting}
\end{Shaded}

The predictions can be made also to future time points of new stations
by using the trend function:

\begin{Shaded}
\begin{Highlighting}[]
\NormalTok{test\_space\_time }\OtherTok{\textless{}{-}} \FunctionTok{subset}\NormalTok{(pollution\_northern\_italy, AirQualityStation }\SpecialCharTok{\%in\%}\NormalTok{ test\_stations }\SpecialCharTok{\&}
\NormalTok{                          time }\SpecialCharTok{\textless{}} \FunctionTok{as.Date}\NormalTok{(}\StringTok{"2022{-}05{-}08"}\NormalTok{) }\SpecialCharTok{\&}\NormalTok{ time }\SpecialCharTok{\textgreater{}} \FunctionTok{as.Date}\NormalTok{(}\StringTok{"2022{-}04{-}30"}\NormalTok{))}
\NormalTok{pred\_space\_time }\OtherTok{\textless{}{-}} \FunctionTok{predict}\NormalTok{(fit, test\_space\_time, }\AttributeTok{type =} \StringTok{"trend"}\NormalTok{,}
                  \AttributeTok{coords =} \FunctionTok{c}\NormalTok{(}\StringTok{"Longitude"}\NormalTok{, }\StringTok{"Latitude"}\NormalTok{), }\AttributeTok{time =} \StringTok{"time"}\NormalTok{,}
                  \AttributeTok{space =} \StringTok{"data"}\NormalTok{)}
\FunctionTok{round}\NormalTok{(}\FunctionTok{c}\NormalTok{(}\AttributeTok{MAE =} \FunctionTok{mean}\NormalTok{(}\FunctionTok{abs}\NormalTok{(pred\_space\_time}\SpecialCharTok{$}\NormalTok{mean\_O3 }\SpecialCharTok{{-}}\NormalTok{ test\_space\_time}\SpecialCharTok{$}\NormalTok{mean\_O3)),}
        \AttributeTok{sd =} \FunctionTok{sd}\NormalTok{(test\_space\_time}\SpecialCharTok{$}\NormalTok{mean\_O3)), }\DecValTok{1}\NormalTok{)              }
\CommentTok{\#\textgreater{}  MAE   sd }
\CommentTok{\#\textgreater{} 18.2 24.3}
\end{Highlighting}
\end{Shaded}

\subsection{Interpretation}\label{interpretation}

\texttt{mshap()} measures how much each latent component contributes to
each observed variable by mean absolute kernel SHAP values of the
decoder. The scaled values sum to one for each variable. Sipilä et
al.~(2026, \emph{Spatial Statistics}) used them to interpret the latent
components of weather and air pollution data over Northern Italy.

\begin{Shaded}
\begin{Highlighting}[]
\NormalTok{ms }\OtherTok{\textless{}{-}} \FunctionTok{mshap}\NormalTok{(fit, }\AttributeTok{n\_sample =} \DecValTok{500}\NormalTok{, }\AttributeTok{n\_background =} \DecValTok{50}\NormalTok{, }\AttributeTok{seed =} \DecValTok{1}\NormalTok{)}
\NormalTok{ms}
\CommentTok{\#\textgreater{} Scaled mean absolute SHAP values (rows: variables, columns: components)}
\CommentTok{\#\textgreater{} }
\CommentTok{\#\textgreater{}             IC1   IC2   IC3}
\CommentTok{\#\textgreater{} mean\_O3   0.364 0.584 0.053}
\CommentTok{\#\textgreater{} mean\_NO2  0.815 0.096 0.089}
\CommentTok{\#\textgreater{} mean\_PM10 0.375 0.111 0.514}
\CommentTok{\#\textgreater{} Average   0.518 0.264 0.219}
\CommentTok{\#\textgreater{} }
\CommentTok{\#\textgreater{} Based on 500 observations and 50 background observations.}
\end{Highlighting}
\end{Shaded}

\begin{Shaded}
\begin{Highlighting}[]
\FunctionTok{plot}\NormalTok{(ms)}
\end{Highlighting}
\end{Shaded}

\includegraphics[width=1\linewidth]{figures/mshap-1}

\subsection{Spatial data formats}\label{spatial-data-formats}

With an \texttt{sf} or \texttt{sftime} object the coordinates and time
points are taken from the object, and the results are again \texttt{sf}
objects. For points, the coordinates are those of the geometry:

\begin{Shaded}
\begin{Highlighting}[]
\NormalTok{s }\OtherTok{\textless{}{-}}\NormalTok{ sf}\SpecialCharTok{::}\FunctionTok{st\_as\_sf}\NormalTok{(dat, }\AttributeTok{coords =} \FunctionTok{c}\NormalTok{(}\StringTok{"lon"}\NormalTok{, }\StringTok{"lat"}\NormalTok{))}
\NormalTok{fit\_sf }\OtherTok{\textless{}{-}} \FunctionTok{iVAEar}\NormalTok{(s, aux, }\AttributeTok{time =} \StringTok{"time"}\NormalTok{, }\AttributeTok{vars =}\NormalTok{ vars\_st, }\AttributeTok{latent\_dim =} \DecValTok{3}\NormalTok{,}
                 \AttributeTok{epochs =} \DecValTok{60}\NormalTok{, }\AttributeTok{verbose =} \DecValTok{0}\NormalTok{, }\AttributeTok{seed =} \DecValTok{1}\NormalTok{)}
\FunctionTok{head}\NormalTok{(}\FunctionTok{components}\NormalTok{(fit\_sf), }\DecValTok{3}\NormalTok{)}
\CommentTok{\#\textgreater{} Simple feature collection with 3 features and 4 fields}
\CommentTok{\#\textgreater{} Geometry type: POINT}
\CommentTok{\#\textgreater{} Dimension:     XY}
\CommentTok{\#\textgreater{} Bounding box:  xmin: 0.2655087 ymin: 0.2936034 xmax: 0.5728534 ymax: 0.9128759}
\CommentTok{\#\textgreater{} CRS:           NA}
\CommentTok{\#\textgreater{}   time        IC1       IC2        IC3                    geometry}
\CommentTok{\#\textgreater{} 1    1 {-}0.1744864 0.5799103 {-}1.0028925 POINT (0.2655087 0.9128759)}
\CommentTok{\#\textgreater{} 2    1  1.0015271 0.3588701 {-}0.1335403 POINT (0.3721239 0.2936034)}
\CommentTok{\#\textgreater{} 3    1  0.5771135 0.6203762 {-}1.1136303 POINT (0.5728534 0.4590657)}
\end{Highlighting}
\end{Shaded}

Spatio-temporal objects of the \textbf{sp} and \textbf{spacetime}
packages (\texttt{STFDF}) are also accepted, and the results are
returned as \texttt{STFDF} objects:

\begin{Shaded}
\begin{Highlighting}[]
\NormalTok{sp\_pts }\OtherTok{\textless{}{-}}\NormalTok{ sp}\SpecialCharTok{::}\FunctionTok{SpatialPoints}\NormalTok{(xy)}
\NormalTok{st }\OtherTok{\textless{}{-}}\NormalTok{ spacetime}\SpecialCharTok{::}\FunctionTok{STFDF}\NormalTok{(sp\_pts, }\FunctionTok{as.POSIXct}\NormalTok{(}\StringTok{"2020{-}01{-}01"}\NormalTok{, }\AttributeTok{tz =} \StringTok{"UTC"}\NormalTok{) }\SpecialCharTok{+} \DecValTok{86400} \SpecialCharTok{*}\NormalTok{ (}\FunctionTok{seq\_len}\NormalTok{(n\_t) }\SpecialCharTok{{-}} \DecValTok{1}\NormalTok{),}
                       \FunctionTok{data.frame}\NormalTok{(x\_st))}
\NormalTok{fit\_stfdf }\OtherTok{\textless{}{-}} \FunctionTok{iVAEar}\NormalTok{(st, aux, }\AttributeTok{latent\_dim =} \DecValTok{3}\NormalTok{, }\AttributeTok{epochs =} \DecValTok{20}\NormalTok{, }\AttributeTok{verbose =} \DecValTok{0}\NormalTok{, }\AttributeTok{seed =} \DecValTok{1}\NormalTok{)}
\FunctionTok{class}\NormalTok{(}\FunctionTok{components}\NormalTok{(fit\_stfdf))}
\end{Highlighting}
\end{Shaded}

\subsection{Large data}\label{large-data}

For large spatio-temporal data the auxiliary matrix can become very
large. With \texttt{batch\_aux\ =\ TRUE} the auxiliary variables are
computed for each batch inside the network from the coordinates and time
points, which gives the same model with much less memory:

\begin{Shaded}
\begin{Highlighting}[]
\NormalTok{fit }\OtherTok{\textless{}{-}} \FunctionTok{iVAEar}\NormalTok{(pollution\_northern\_italy, }\FunctionTok{aux\_rbf}\NormalTok{(),}
              \AttributeTok{coords =} \FunctionTok{c}\NormalTok{(}\StringTok{"Longitude"}\NormalTok{, }\StringTok{"Latitude"}\NormalTok{), }\AttributeTok{time =} \StringTok{"time"}\NormalTok{, }\AttributeTok{vars =}\NormalTok{ vars,}
              \AttributeTok{batch\_aux =} \ConstantTok{TRUE}\NormalTok{)}
\end{Highlighting}
\end{Shaded}

\subsection{Saving models}\label{saving-models}

The trained weights are stored in the fitted object, so models can be
saved with \texttt{save\_nlbss()} (or \texttt{saveRDS()}) and used in a
new session:

\begin{Shaded}
\begin{Highlighting}[]
\NormalTok{file }\OtherTok{\textless{}{-}} \FunctionTok{tempfile}\NormalTok{(}\AttributeTok{fileext =} \StringTok{".rds"}\NormalTok{)}
\FunctionTok{save\_nlbss}\NormalTok{(fit, file)}
\NormalTok{fit2 }\OtherTok{\textless{}{-}} \FunctionTok{load\_nlbss}\NormalTok{(file)}
\FunctionTok{all.equal}\NormalTok{(}\FunctionTok{predict}\NormalTok{(fit2, }\AttributeTok{type =} \StringTok{"trend"}\NormalTok{), }\FunctionTok{predict}\NormalTok{(fit, }\AttributeTok{type =} \StringTok{"trend"}\NormalTok{))}
\CommentTok{\#\textgreater{} [1] TRUE}
\end{Highlighting}
\end{Shaded}

\subsection{References}\label{references}

\begin{itemize}
\tightlist
\item
  Hyvärinen, A. and Morioka, H. (2016). Unsupervised feature extraction
  by time-contrastive learning and nonlinear ICA. \emph{Advances in
  Neural Information Processing Systems}, 29.
\item
  Khemakhem, I., Kingma, D., Monti, R. and Hyvärinen, A. (2020).
  Variational autoencoders and nonlinear ICA: A unifying framework.
  \emph{AISTATS 2020}, 2207--2217.
\item
  Sipilä, M., Nordhausen, K. and Taskinen, S. (2024). Nonlinear blind
  source separation exploiting spatial nonstationarity.
  \emph{Information Sciences}, 665,

  \begin{enumerate}
  \def\labelenumi{\arabic{enumi}.}
  \setcounter{enumi}{120364}
  \tightlist
  \item
    \url{https://doi.org/10.1016/j.ins.2024.120365}
  \end{enumerate}
\item
  Sipilä, M., Cappello, C., De Iaco, S., Nordhausen, K. and Taskinen, S.
  (2025). Modelling multivariate spatio-temporal data with identifiable
  variational autoencoders. \emph{Neural Networks}, 181, 106774.
  \url{https://doi.org/10.1016/j.neunet.2024.106774}
\item
  Sipilä, M., Nordhausen, K. and Taskinen, S. (2025). Identifiable
  autoregressive variational autoencoders for nonlinear and
  nonstationary spatio-temporal blind source separation. \emph{ECML PKDD
  2025}, Lecture Notes in Computer Science 16019, 362--380.
  \url{https://doi.org/10.1007/978-3-032-06109-6_21}
\item
  Sipilä, M., De Iaco, S., Cappello, C., Nordhausen, K., Palma, M. and
  Taskinen, S. (2026). Enhancing identifiable variational autoencoder in
  the presence of missing values and auxiliary covariates for
  spatio-temporal ozone modeling. \emph{Environmetrics}, 37(6), e70121.
  \url{https://doi.org/10.1002/env.70121}
\item
  Sipilä, M., Cappello, C., De Iaco, S., Nordhausen, K., Palma, M. and
  Taskinen, S. (2026). Identifiable variational autoencoders for
  interpretable dimension reduction and prediction in multivariate
  spatio-temporal processes. \emph{Spatial Statistics}, 101054.
  \url{https://doi.org/10.1016/j.spasta.2026.101054}
\end{itemize}

\end{document}
